\SourceLeaf{186}
\SeventhArticle{11}
Quoties $a, a', a''$ etc. vel ex parte vel omnes sunt irrationales,. utilis erit transformatio functionis fractae, ex qua numeros $R, R', R''$ etc. derivamus, in functionem integram: principia talis transformationis, quum in libris algebraicis non inveniantur, hoc loco breviter explicabimus. Propositis scilicet tribus functionibus integris $Z,\zeta,\zeta'$ indeterminatae $z$, quaeritur functio integra, quae fractae $\frac Z\zeta$ vice fungi possit, quatenus pro $z$ accipitur radix quaecunque aequationis $\zeta'=0$. Supponemus autem, $\zeta$ pro nullo horum valorum ipsius $z$ evanescere, sive quod eodem redit, $\zeta$ atque $\zeta'$. nullum divisorem communem indeterminatum implicare. Exponentes potestatum altissimarum ipsius $z$ in $\zeta$ atque $\zeta'$ per $k,k'$ denotabimus.

Dividatur sueto more $\zeta$ per $\zeta'$, donec residui ordo infra $k'$ depressus sit; statuatur residuum $=\frac{\zeta''}{\lambda}$, eiusque ordo $=k''$, ita ut $\frac1\lambda z^{k''}$ sit residui terminus altissimus; divisionis quotientem ponemus $=\frac p\lambda$. Perinde ex divisione functionis $\zeta'$ per $\zeta''$ prodeat residuum $\frac{\zeta'''}{\lambda'}$ ordinis $k'''$, quotiens $\frac{p'}{\lambda'}$; dein rursus e divisione functionis $\zeta''$ per $\zeta'''$ prodeat residuum $\frac{\zeta''''}{\lambda''}$ ordinis $k''''$ atque quotiens $\frac{p''}{\lambda''}$ et sic porro, donec in serie functionum $\zeta'',\ \zeta'''.\ \zeta''''$ etc., quae singulae terminum suum altissimum co\"efficiente $1$ affectum habebunt, perveniatur ad $\zeta^{(m)}=1$. Hoc tandem evenire debere facile perspicitur, quum quaelibet functionum $\zeta,\ \zeta'\ \zeta'',\ \zeta'''$ etc. cum praecedenti divisorem communem indeterminatum habere nequeat, adeoque certo divisio absque residuo fieri nequeat, quamdiu divisor fuerit ordinis maioris quam $0$. Habebimus igitur seriem aequationum
\[
\begin{aligned}
\zeta''&=\lambda\zeta-p\zeta'\\
\zeta'''&=\lambda'\zeta'-p'\zeta''\\
\zeta''''&=\lambda''\zeta''-p''\zeta'''\\
\zeta'''''&=\lambda'''\zeta'''-p'''\zeta''''\\
&\text{etc. usque ad}\\
\zeta^{(m)}&=\lambda^{(m-2)}\zeta^{(m-2)}-p^{(m-2)}\zeta^{(m-1)}
\end{aligned}
\]
ubi $\zeta'',\allowbreak \ \zeta''',\allowbreak \ \zeta''''\ldotp\ldotp\ldotp\ldotp\zeta^{(m)}$ sunt functiones integrae ipsius $z$ ordinis $k'',\allowbreak \ k''',\allowbreak \ k''''\ldotp\ldotp\ldotp\ldotp k^{(m)}$; numeri $k',\allowbreak \ k'',\allowbreak \ k'''\ldotp\ldotp\ldotp\ldotp k^{(m)}$ continuo decrescentes usque ad ultimum $k^{(m)}=0$; $p,\allowbreak p',\allowbreak p'',\allowbreak p'''$ etc. quoque functiones integrae ipsius $z$ ordinis $k-k',\allowbreak \ k'-k'',\allowbreak \ k''-k''',\allowbreak \ k'''-k''''$ etc. (excepto casu, ubi $k<k'$, in quo manifesto statui debet $p=0$).

His ita praeparatis formamus \textit{secundam} seriem functionum integrarum ipsius $z$, puta $\eta,\ \eta',\ \eta'',\ \eta'''\ldotp\ldotp\ldotp\ldotp\eta^{(m)}$. Et quidem statuemus $\eta=1,\ \eta'=0$, reliquas vero
\SourceLeaf{187}
singulas e binis praecedentibus per eandem legem derivamus, per quam functiones $\zeta,\zeta',\zeta'',\zeta'''$ etc. inter se nexae sunt, scilicet per aequationes
\[
\begin{aligned}
\eta''&=\lambda\eta-p\eta'\\
\eta'''&=\lambda'\eta'-p'\eta''\\
\eta''''&=\lambda''\eta''-p''\eta'''\\
\eta'''''&=\lambda'''\eta'''-p'''\eta''''\quad\text{etc. usque ad}\\
\eta^{(m)}&=\lambda^{(m-2)}\eta^{(m-2)}-p^{(m-2)}\eta^{(m-1)}
\end{aligned}
\]
Manifesto $\eta''=\lambda$ hic est ordinis $0$; $\eta'''=-\lambda p'$ ordinis $k'-k''$, et perinde sequentes $\eta'''',\eta'''''$ etc. resp. ordinis $k'-k''',\ k'-k''''$ etc., ita ut ultima $\eta^{(m)}$ ascendat ad ordinem $k'-k^{(m-1)}$.

Porro consideremus \textit{tertiam} functionum seriem, $\zeta-\zeta\eta,\ \zeta'-\zeta\eta',\ \zeta''-\zeta\eta'',\ \zeta'''-\zeta\eta'''$ etc., inter cuius terminos quosvis ternos consequentes manifesto similis relatio intercedet, scilicet
\[
\begin{aligned}
\zeta''-\zeta\eta''&=\lambda(\zeta-\zeta\eta)-p(\zeta'-\zeta\eta')\\
\zeta'''-\zeta\eta'''&=\lambda'(\zeta'-\zeta\eta')-p'(\zeta''-\zeta\eta'')\\
\zeta''''-\zeta\eta''''&=\lambda''(\zeta''-\zeta\eta'')-p''(\zeta'''-\zeta\eta''')
\end{aligned}
\]
Iam prima harum functionum fit $=0$, secunda $=\zeta'$: hinc facile colligitur, singulas per $\zeta'$ divisibiles fore.

Hinc autem nullo negotio sequitur, loco fractionis $\frac Z\zeta$ adoptari posse functionem integram $Z\eta^{(m)}$, quatenus quidem ipsi $z$ non tribuantur alii valores nisi qui sint radices aequationis $\zeta'=0$: manifesto enim differentia $\frac{Z(1-\zeta\eta^{(m)})}{\zeta}$ pro tali valore ipsius $z$ necessario evanescit, quum $1-\zeta\eta^{(m)}=\zeta^{(m)}-\zeta\eta^{(m)}$ per $\zeta'$ sit divisibilis.

Loco functionis $Z\eta^{(m)}$ etiam adoptari poterit eius residuum ex divisione per $\zeta'$ ortum, cuius ordo erit inferior ordine functionis $\zeta'$.

Ceterum hocce residuum commodius per algorithmum sequentem immediate eruere licet. Formentur aequationes sequentes
\[
\begin{aligned}
Z&=q'\zeta'+Z'\\
Z'&=q''\zeta''+Z''\\
Z''&=q'''\zeta'''+Z'''\\
Z'''&=q''''\zeta''''+Z''''\quad\text{etc. usque ad}\\
Z^{(m-1)}&=q^{(m)}\zeta^{(m)}+Z^{(m)}
\end{aligned}
\]
\SourceLeaf{188}
scilicet deinceps dividendo $Z$ per $\zeta'$, dein residuum primae divisionis $Z'$ per $\zeta''$, tum residuum secundae divisionis per $\zeta'''$ ac sic porro. Quum residuum semper ad ordinem inferiorem pertineat quam divisor, ordo functionum $Z',Z'',Z''',Z''''$ etc. erit resp. inferior quam $k',k'',k''',k''''$ etc.; ultima vero $Z^{(m)}$ necessario fit $=0$, quum divisor $\zeta^{(m)}$ sit $=1$. Habemus itaque
\[
Z=q'\zeta'+q''\zeta''+q'''\zeta'''+q''''\zeta''''+\text{etc.}+q^{(m)}\zeta^{(m)}
\]
Quatenus autem pro $z$ solae radices aequationis $\zeta'=0$ accipiuntur, fit $\zeta'=0$, $\zeta''=\zeta\eta''$, $\zeta'''=\zeta\eta'''$, $\zeta''''=\zeta\eta''''$ etc., unde sub eadem restrictione erit
\[
\frac{Z}{\zeta}=q''\eta''+q'''\eta'''+q''''\eta''''+\text{etc.}+q^{(m)}\eta^{(m)}
\]
Ordo vero huius expressionis necessario erit infra $k'$: quum enim ordo quotientium $q'',q''',q''''$ etc. esse debeat infra $k'-k''$, $k''-k'''$, $k'''-k''''$ etc., ordo singularum partium $q''\eta''$, $q'''\eta'''$, $q''''\eta''''$ etc. erit infra $k'-k''$, $k'-k'''$, $k'-k''''$ etc.

Denique adhuc observamus, si forte inter valores indeterminatae $z$, quos in fractione $\frac{Z}{\zeta}$ substituere oporteat, rationales cum irrationalibus mixti reperiantur, magis e re fore, illos ab his separare atque hos solos in aequatione $\zeta'=0$ comprehendere. Pro rationalibus enim valoribus calculi compendio opus non erit; pro irrationalibus autem calculus tanto simplicior erit, quo minor fuerit gradus functionis integrae, ad quam fractam reducere licet.

\SeventhArticle{12}
Ecce nunc exemplum transformationis in art. praec explicatae. Proposita sit functio fracta
\[
\frac{z^6-\frac{50}{39}z^4+\frac{283}{715}zz-\frac{256}{15015}}{7z^6-\frac{105}{13}z^4+\frac{315}{143}zz-\frac{35}{429}}
\]
in qua $z$ indefinite repraesentat radices aequationis
\[
z^7-\frac{21}{13}z^5+\frac{105}{143}z^3-\frac{35}{429}z=0
\]
Si hic omnes septem radices complecti vellemus, ad functionem integram sexti ordinis delaberemur. Manifesto autem pro valore rationali $z=0$ computus fractionis obvius est, datque valorem $\frac{256}{1225}$: quapropter seposita hac radice in aequatione sexti gradus subsistemus:
\SourceLeaf{189}
\[
z^6-\frac{21}{13}z^4+\frac{105}{143}zz-\frac{35}{429}=0
\]
quo pacto facile praevidemus orturam esse functionem integram quarti ordinis. Iam ex applicatione praeceptorum praecedentium prodeunt sequentia:
\[
\begin{aligned}
\zeta&=7z^6-\frac{105}{13}z^4+\frac{315}{143}zz-\frac{35}{429}\\[5pt]
\zeta'&=z^6-\frac{21}{13}z^4+\frac{105}{143}zz-\frac{35}{429}\\[5pt]
\zeta''&=z^4-\frac{10}{11}zz+\frac{5}{33}\\[5pt]
\zeta'''&=zz-\frac{3}{7}\\[5pt]
\zeta''''&=1
\end{aligned}\]
\[\begin{aligned}
\lambda&=\frac{13}{42}&p&=\frac{13}{6}\\[5pt]
\lambda'&=-\frac{4719}{280}&p'&=-\frac{4719}{280}zz+\frac{3333}{280}\\[5pt]
\lambda''&=-\frac{147}{8}&p''&=-\frac{147}{8}zz+\frac{777}{88}
\end{aligned}\]
\[\begin{aligned}
\eta&=1\\[5pt]
\eta'&=0\\[5pt]
\eta''&=\frac{13}{42}\\[5pt]
\eta'''&=\frac{20449}{3920}zz-\frac{14443}{3920}\\[5pt]
\eta''''&=\frac{61347}{640}z^4-\frac{127413}{1120}zz+\frac{120263}{4480}
\end{aligned}\]
\[\begin{aligned}
Z&=z^6-\frac{50}{39}z^4+\frac{283}{715}zz-\frac{256}{15015};&q'&=1\\[5pt]
Z'&=\frac{1}{3}z^4-\frac{22}{65}zz+\frac{323}{5005}&q''&=\frac{1}{3}\\[5pt]
Z''&=-\frac{76}{2145}zz+\frac{632}{45045}&q'''&=-\frac{76}{2145}\\[5pt]
Z'''&=-\frac{4}{3465}&q''''&=-\frac{4}{3465}
\end{aligned}
\]
Hinc tandem derivatur functio integra fractioni propositae aequivalens:
\[
-\frac{1859}{16800}z^4-\frac{1573}{29400}zz+\frac{7947}{39200}
\]

\SeventhArticle{13}
Ad determinandum gradum praecisionis, qua formula nostra integralis $RA+R'A'+R''A''+\text{etc.}+R^{(n)}A^{(n)}$ gaudet, statuamus generaliter
\[
Ra^m+R'a'{}^m+R''a''{}^m+\text{etc.}+R^{(n)}a^{(n)m}=\frac{1}{m+1}-k^{(m)}
\]
ita ut $k^{(m)}$ sit differentia inter integralis $\int t^m\,\mathrm{d}t$ a $t=0$ usque ad $t=1$ sumti valorem verum atque approximatum. Habebimus itaque, singulis fractionibus in series evolutis,
\SourceLeaf{190}
\[
\begin{aligned}
&\frac{R}{t-a}+\frac{R'}{t-a'}+\frac{R''}{t-a''}+\text{etc.}+\frac{R^{(n)}}{t-a^{(n)}}\\
&\quad=(1-k)t^{-1}+(\tfrac12-k')t^{-2}+(\tfrac13-k'')t^{-3}\\ &\qquad+(\tfrac14-k''')t^{-4}+\text{etc.}\\
&\quad=t^{-1}+\tfrac12t^{-2}+\tfrac13t^{-3}+\tfrac14t^{-4}+\text{etc.}-\theta
\end{aligned}
\]
si statuimus
\[
\theta=kt^{-1}+k't^{-2}+k''t^{-3}+k'''t^{-4}+\text{etc.}
\]
sive potius (quum iam sciamus, $k$, $k'$, $k''$, $k'''$ etc. usque ad $k^{(n)}$ sponte evanescere debere)
\[
\theta=k^{(n+1)}t^{-(n+2)}+k^{(n+2)}t^{-(n+3)}+k^{(n+3)}t^{-(n+4)}+\text{etc.}
\]

Multiplicando per $T$ fit
\[
\begin{aligned}&T\left(\frac{R}{t-a}+\frac{R'}{t-a'}+\frac{R''}{t-a''}+\text{etc.}+\frac{R^{(n)}}{t-a^{(n)}}\right)\\ &=T'+T''-T\theta\end{aligned}
\]
Pars prior huius aequationis est functio integra ipsius $t$ ordinis $n$, eiusque valores determinati pro $t=a$, $t=a'$, $t=a''$ etc. resp. fiunt $MR$, $M'R'$, $M''R''$ etc.: quapropter, quum eadem valeant de functione $T'$, uti ex ipso modo numeros $R$, $R'$, $R''$ etc. determinandi perspicuum est, necessario illa pars prior aequationis identica esse debet cum $T'$, adeoque $T''=T\theta$. Oritur itaque $\theta$ ex evolutione fractionis $\frac{T''}{T}$, quo pacto co\"efficientes $k^{(n+1)}$, $k^{(n+2)}$ etc. quousque libet determinari poterunt. Quibus inventis correctio valoris nostri approximati integralis $\int y\,\mathrm{d}t$ erit
\[
=k^{(n+1)}K^{(n+1)}+k^{(n+2)}K^{(n+2)}+\text{etc.}
\]
si series, in quam evolvitur $y$, est
\[
y=K+K't+K''tt+K'''t^3+\text{etc.}
\]

\SeventhArticle{14}
Si magis placet, correctionem exprimere per co\"efficientes seriei secundum potestates ipsius $t-\tfrac12$ progredientis
\[
y=L+L'(t-\tfrac12)+L''(t-\tfrac12)^2+L'''(t-\tfrac12)^3+\text{etc.}
\]
illa erit
\SourceLeaf{191}
\[
=l^{(n+1)}L^{(n+1)}+l^{(n+2)}L^{(n+2)}+l^{(n+3)}L^{(n+3)}+\text{etc.}
\]
si generaliter per $l^{(m)}$ exprimimus correctionem valoris approximati integralis $\int(t-\tfrac12)^m\,\mathrm{d}t$. Hae correctiones $l^{(m)}$ cum correctionibus $k^{(m)}$ nexae erunt per aequationem
\[
\begin{aligned}&l^{(m)}=k^{(m)}-\tfrac12mk^{(m-1)}+\tfrac14\cdot\frac{m\cdot m-1}{1\cdot2}k^{(m-2)}\\ &-\tfrac18\cdot\frac{m\cdot m-1\cdot m-2}{1\cdot2\cdot3}k^{(m-3)}+\text{etc.}\end{aligned}
\]
Quo vero illas independenter eruere possimus, perpendamus, functionem $\theta$ per substitutionem $t=\tfrac12u+\tfrac12$ transire in
\[
\begin{aligned}
&\phantom{{}+{}}2k(u^{-1}-u^{-2}+u^{-3}-u^{-4}+\text{etc.})\\
&+4k'(u^{-2}-2u^{-3}+3u^{-4}-4u^{-5}+\text{etc.})\\
&+8k''(u^{-3}-3u^{-4}+6u^{-5}-10u^{-6}+\text{etc.})\\
&+16k'''(u^{-4}-4u^{-5}+10u^{-6}-20u^{-7}+\text{etc.})\\
&+\text{etc.}
\end{aligned}
\]
sive in
\[
\begin{aligned}
&2ku^{-1}+4(k'-\tfrac12)u^{-2}+8(k''-\tfrac12\cdot2k'+\tfrac14k)u^{-3}\\
&\qquad+16(k'''-\tfrac12\cdot3k''+\tfrac14\cdot3k'-\tfrac18k)u^{-4}+\text{etc.}
\end{aligned}
\]
sive in
\[
2lu^{-1}+4l'u^{-2}+8l''u^{-3}+16l'''u^{-4}+\text{etc.}
\]
sive denique, quum a priori sciamus, $l$, $l'$, $l''$, $l'''$ etc. usque ad $l^{(n)}$ sponte evanescere, in
\[
\begin{aligned}&2^{n+2}l^{(n+1)}u^{-(n+2)}+2^{n+3}l^{(n+2)}u^{-(n+3)}\\ &+2^{n+4}l^{(n+4)}u^{-(n+4)}+\text{etc.}\end{aligned}
\]
At $\theta=\frac{T''}{T}$; quare quum $T$, $T''$ per substitutionem $t=\tfrac12u+\tfrac12$ transeant in $\frac{U}{2^{n+1}}$, $\frac{U''}{2^n}$, (art. 9), functio $\theta$ per eandem substitutionem transibit in $\frac{2U''}{U}$. Quodsi itaque seriem ex evolutione fractionis $\frac{U''}{U}$ oriundam per $\Omega$ designamus, erit
\[
\begin{aligned}&\Omega=2^{n+1}l^{(n+1)}u^{-(n+2)}+2^{n+2}l^{(n+2)}u^{-(n+3)}\\ &+2^{n+3}l^{(n+3)}u^{-(n+4)}+\text{etc.}\end{aligned}
\]
quo pacto co\"efficientes $l^{(n+1)}$, $l^{(n+2)}$ etc. quousque lubet erui poterunt.

Ita in exemplo art. 10 invenimus
\[
\begin{aligned}
U''&=-\tfrac{176}{13125}u^{-1}-\tfrac{304}{28125}u^{-3}-\tfrac{2576}{309375}u^{-5}-\text{etc.}\\
\Omega&=-\tfrac{176}{13125}u^{-7}-\tfrac{832}{28125}u^{-9}-\tfrac{189856}{4296875}u^{-11}-\text{etc.}
\end{aligned}
\]
